\section{Classical Hamiltonian observables on the torus family}
\label{hp:sec}

A compact complex torus has finite-dimensional Dolbeault cohomology,
but replacing differential forms by those cohomology groups does not by itself
compactify a field theory. The differential must still describe the
variation of the fiber, and the pairing must still give the bracket
of observables. In a presymplectic theory there is a further issue:
Hamiltonian observables depend on a resolution of the kernel of the
presymplectic operator. A comparison of field complexes need not
compare these resolutions.

We address these questions for the classical observable model of the
alternative holomorphic twist of the abelian six-dimensional
$(2,0)$ theory. Fix the sufficiently small nonzero parameter $u$,
integral lifts and adapted filling of
Theorem~\ref{cmp:thm:comparison}. Put
\[
 U=\mathbb P^1\setminus\{0,1,\infty\},\qquad
 X=Y^{\mathrm{ad}}(u)|_U,\qquad f:X\longrightarrow U.
\]
Thus $f$ is the actual smooth proper family of complex two-tori.
We construct a complex of finite-rank smooth bundles on $U$ whose
compactly supported sections give the same resolved classical
observables as forms on $X$. Its differential contains the derivative
of the full period matrix. The inclusion of invariant forms preserves
the Poisson bracket, while averaging supplies a chain inverse.
The free finite covers at $0$ and $1$ have a separate descent formula
with an essential degree normalization. The singular fiber at infinity
requires an additional comparison.

\subsection{The resolved observable complex}
\label{hp:subsec:observables}

All coefficients in this section are complex. A $P_0$ algebra is a
commutative differential graded algebra with a Poisson bracket of
cohomological degree $+1$. For bidegrees $(i,\epsilon)$ and
$(j,\epsilon')$ in $\mathbb Z\times\mathbb Z/2$, the symmetry sign is
$(-1)^{ij+\epsilon\epsilon'}$. The alternative regrading used here is
concentrated in even supplemental $\mathbb Z/2$ degree
\cite[\S4.5.1]{SaberiWilliamsConstraints}. Thus $\epsilon=0$ for all
the generators below; their cohomological degrees still determine the
Koszul signs.

For an open set $V$ in a complex threefold, let
$\Omega_c^{p,q}(V)$ be its smooth compactly supported $(p,q)$-forms.
We use cohomological shifts, so an element of degree $d$ has degree
$d-r$ in a complex shifted by $[r]$. Define
\begin{equation}\label{hp:eq:linear}
\begin{split}
 V_H(V)&=V_\chi(V)\oplus V_\beta(V)\oplus V_\gamma(V),\\
 V_\chi(V)&=\Omega_c^{\leq1,\bullet}(V)[3],\qquad
 V_\beta(V)=\Omega_c^{3,\bullet}(V)[3],\qquad
 V_\gamma(V)=\Omega_c^{0,\bullet}(V)[1].
\end{split}
\end{equation}
Here $\Omega^{\leq1,\bullet}$ is the total complex of columns
$p=0,1$. The degrees of its $(p,q)$-forms are $p+q-3$ after the
shift. The degrees in $V_\beta$ and $V_\gamma$ are $q-3$ and $q-1$,
respectively. The linear differential is
\begin{equation}\label{hp:eq:differential}
 D_H|_{V_\chi}=\bar\partial+\partial_{\leq1},\qquad
 D_H|_{V_\beta\oplus V_\gamma}=\bar\partial.
\end{equation}
The truncated operator $\partial_{\leq1}$ maps the $p=0$ column
to the $p=1$ column and is zero on the latter. There is no
$\partial$ differential between the $\beta$ and $\gamma$ summands.

Fix a symmetric-tensor realization, denoted by $\mathsf S$.
The algebraic realization is
$\mathsf S(V)=\bigoplus_{n\geq0}(V^{\otimes n})_{\mathfrak S_n}$.
In the degreewise completed realization, replace each finite tensor
power by its completed projective tensor power, with the topology on
compact forms given by the inductive limit of the fixed compact-support
Fr\'echet spaces. In both cases the sum over polynomial degree is a
direct sum, and the $n=0$ summand is $\mathbb C$ with unit $1$.
Every comparison below uses the same realization on its two sides.
No identification of the algebraic and completed realizations, or
replacement of the polynomial direct sum by a product, is intended.

The resolved classical observable algebra is
\begin{equation}\label{hp:eq:observables}
 \mathcal O_H(V)=\bigl(\mathsf S(V_H(V)),D_H\bigr).
\end{equation}
Its nonzero brackets on linear generators are
\begin{equation}\label{hp:eq:brackets}
 B_\chi(\alpha,\alpha')=\int_V\alpha\wedge\partial\alpha',
 \qquad
 B_{\beta\gamma}(\beta,\gamma)=\int_V\beta\wedge\gamma.
\end{equation}
The operator in $B_\chi$ is the full holomorphic exterior derivative
of the representing form. Only the internal differential in
\eqref{hp:eq:differential} is truncated. Integration selects bidegree
$(3,3)$, and the shifts make both brackets have degree $+1$.
The bracket on polynomials is their extension by the graded Leibniz
rule. For disjoint open sets, extend linear generators by zero and
multiply the resulting polynomials. These are the factorization
operations, and the scalar summand supplies their unit.

The Hamiltonian qualification specifies which observable model is
being used. If $\mathcal E$ is a free presymplectic field complex and
$\omega:\mathcal E\to\mathcal E^![-1]$ is its presymplectic
operator, the symbol $!$ means density dual. Choose a semi-free
resolution $\mathcal K_\omega$ of $\ker\omega$, together with its
differential-operator map $\pi_\omega:\mathcal K_\omega\to\mathcal E$.
The smoothed derived invariants have linear generators
\begin{equation}\label{hp:eq:kernel-recipe}
 (\mathcal K_\omega)^!_c(V)\oplus\mathcal E^!_c(V)[1],
\end{equation}
with the differential induced by the two internal differentials and
$\pi_\omega$. This is the kernel-resolution construction of
\cite[\S2.4.1]{SaberiWilliamsConstraints}. Its stated hypotheses,
including the cone-resolution condition for the chiral two-form,
give the explicit model \eqref{hp:eq:linear}--\eqref{hp:eq:brackets}
and its $P_0$ factorization structure
\cite[\S4.6]{SaberiWilliamsConstraints}. We retain that resolved model
and its factorization descent in the chosen tensor convention.
The arguments that follow act directly on its generators. They do
not require an invariance theorem for Hamiltonian observables under
arbitrary quasi-isomorphisms of presymplectic field complexes.

\subsection{Descent through a free finite cover}
\label{hp:subsec:finite}

The finite fillings have ramified maps of base disks. Their maps of
total spaces are nevertheless unramified: the cyclic actions are
free. This distinction determines the descent of compact forms.

Let $\pi:\widetilde X\to X_0$ be a finite holomorphic covering of
degree $m$, given first as a quotient by a free action of a group
$G$ of order $m$. For $V\subset X_0$, write
$\widetilde V=\pi^{-1}V$. Pullback identifies compact forms on $V$
with invariant compact forms on $\widetilde V$. Its inverse is the
normalized sheet trace
\begin{equation}\label{hp:eq:trace}
 r=\frac1m\operatorname{Tr}_\pi,
 \qquad
 \pi^*r=\frac1m\sum_{g\in G}g^*,\qquad r\pi^*=1.
\end{equation}
The trace is defined on an evenly covered open set by adding the
pullbacks along its inverse branches. These local sums agree on
overlaps. Properness of $\pi$ preserves compact support under
pullback, and the image of a compact support is compact under trace.
Both maps commute with $\partial$ and $\bar\partial$.

\begin{proposition}[Normalized finite descent]\label{hp:prop:finite}
On the invariant linear generator complex
$V_H(\widetilde V)^G$, use the bracket
$m^{-1}B_{\widetilde V}$, where $B$ denotes the two blocks in
\eqref{hp:eq:brackets}. Then
\begin{equation}\label{hp:eq:finite-isomorphism}
 \mathsf S(\pi^*):\mathcal O_H(V)
 \xrightarrow{\ \cong\ }
 \left(\mathsf S\bigl(V_H(\widetilde V)^G\bigr),
             D_H,\frac1m B_{\widetilde V}\right)
\end{equation}
is a strict unital $P_0$ isomorphism, with inverse $\mathsf S(r)$.
These isomorphisms commute with all disjoint-open factorization
operations. Equivalently, $m^{-1/2}\pi^*$ identifies the downstairs
algebra with the symmetric algebra on invariant linear generators
inside the original, unrescaled cover algebra.
\end{proposition}

\begin{proof}
For any compactly supported top form $\eta$, holomorphic change of
variables gives
$\int_{\widetilde V}\pi^*\eta=m\int_V\eta$.
Consequently
\[
 B_{\widetilde V}(\pi^*a,\pi^*b)=m B_V(a,b)
\]
for both bracket blocks. This proves bracket compatibility after
division by $m$. Equation~\eqref{hp:eq:trace} gives inverse chain
maps on the specified linear spaces, so the same symmetric functor
gives inverse algebra maps. These maps are continuous on each
fixed compact-support stage and hence in the completed realization
as well. The graded Leibniz rule proves bracket compatibility in
every polynomial degree. The degree-zero maps fix $1$.

Inverse images of disjoint open sets are disjoint. Pullback and
trace commute with extension by zero, so the maps commute with
every factorization operation. Finally, multiplying each linear
pullback by $m^{-1/2}$ cancels the factor $m$ in its bilinear
bracket. This gives the equivalent normalization with the original
cover bracket.
\end{proof}

The target of \eqref{hp:eq:finite-isomorphism} is formed by taking
invariant \emph{linear} generators before applying $\mathsf S$.
It is generally smaller than the invariant algebra
$\mathsf S(V_H(\widetilde V))^G$. To see this already in polynomial
degree two, choose a cohomologically even generator on an evenly covered open set
and denote its cyclic sheet translates by $x_1,\ldots,x_m$.
The invariant linear span is generated by $\sum_i x_i$, whereas
$\sum_i x_i^2$ is not a scalar multiple of $(\sum_i x_i)^2$.
At $(x_1,x_2,x_3,\ldots)=(1,-1,0,\ldots)$ their values are $2$
and $0$. The invariant quadratic space has dimension two for
$m=3$ and three for $m=4$, according to the cyclic distances of
the two indices; the symmetric square of the invariant linear
space has dimension one. This distinguishes the underlying
polynomial carriers. It makes no assertion about the cohomology
of the extra invariant polynomials.

There is also a distinction between the inverse on invariant
generators and averaging from the whole cover. Choose compact
$\beta_i,\gamma_i$ on two disjoint sheets, with
$\int\beta_i\wedge\gamma_i=1$, obtained by translation from the
same downstairs pair. The averages of
$\beta_1-\beta_2$ and $\gamma_1-\gamma_2$ vanish, while their
bracket is $2$. Thus averaging on all cover generators is not a
Poisson retraction. The inverse in
Proposition~\ref{hp:prop:finite} has the smaller invariant domain.

For successive finite covers, sheet traces compose and degrees
multiply. Hence their normalized traces compose as well. This
argument uses local sheets and remains valid when the composite
cover is not Galois.

In the two fillings of the selected threefold, the root maps are
$t=s^3$ and $t=1+s^4$. In the respective original deck markings,
the affine terms are $\delta_0/3$ and $-\delta_1/4$; the selected
increments are $e_4/3$ and $-f_4/4$. The invariant quotient
covector vanishes on $\delta_i$ and evaluates the integral
increments as $1$ and $-1$. Proposition~\ref{prop:affine-free}
therefore gives free central actions of orders three and four.
Away from the central fiber, nontrivial root rotations have no
fixed base point. The total-space quotient maps are thus free
finite covers, and Proposition~\ref{hp:prop:finite} applies with
the full affine actions and degrees $3,4$.

\subsection{Invariant total forms and their differential}
\label{hp:subsec:invariant}

We now return to $f:X\to U$. On a simply connected disk $D\subset U$,
choose a local marking and the holomorphic origin supplied by the
conormal section. The family has a presentation
\begin{equation}\label{hp:eq:torus}
 X_D=\bigl(D\times\mathbb C^2\bigr)/
             \bigl(\mathbb Z^2+\Omega(z)\mathbb Z^2\bigr),
 \qquad \det(\Omega-\bar\Omega)\ne0.
\end{equation}
The period matrix here belongs to the actual divisor-line quotient.
If $\tau$ is the normalized elliptic period and $\mu$ is the
Abel--Jacobi coordinate of the section $P$, the theta-ratio frame
of $\mathcal O(P-O)$ has multiplier $e^{2\pi i\mu}$ under the
$\tau$ period. The further quotient covers translation by $-6P$.
Its multiplier in this frame is independent of the elliptic
coordinate, since the ratio of two such lifts is holomorphic on a
compact elliptic curve. Write that multiplier as $e^{2\pi i\sigma}$.
The four periods are therefore
\begin{equation}\label{hp:eq:native-periods}
 (1,0),\quad(0,1),\quad(\tau,\mu),\quad(-6\mu,\sigma),
 \qquad
 \Omega=\begin{pmatrix}\tau&-6\mu\\ \mu&\sigma\end{pmatrix}.
\end{equation}
The function $\sigma$ includes the chosen lift and the fixed
parameter $u$. It is not assumed constant or proportional to $\tau$.

Write $w=x+\Omega(z)y$, with $x,y$ real modulo $\mathbb Z^2$.
For constant real vectors $a,b$, the translation
$w\mapsto w+a+\Omega(z)b$ is holomorphic on the total space.
These translations form a locally flat real four-torus. Their
normalized Haar average on total forms is denoted by $R$.
A change of integral marking is an automorphism of this compact
torus and preserves Haar measure; changing the origin is a
translation. Thus $R$ is intrinsic.

The relevant invariant forms are absolute forms on $X_D$. They
include forms in the base direction. Put
\begin{equation}\label{hp:eq:coframe}
 \mathsf Y=\Omega-\bar\Omega,\qquad
 \theta=dw-d\Omega\,y=dx+\Omega\,dy,\qquad
 \bar\theta=d\bar w-d\bar\Omega\,y=dx+\bar\Omega\,dy.
\end{equation}
The forms $dz,\theta_1,\theta_2$ and
$d\bar z,\bar\theta_1,\bar\theta_2$ give an invariant
$(1,0)$ and $(0,1)$ coframe. An invariant form is uniquely a
finite sum of their wedge monomials with smooth base coefficients.
Consequently invariant $(p,q)$-forms define a finite-rank smooth
bundle on $U$. Denote its compactly supported sections by
$W^{p,q}(D)$, and set
\begin{equation}\label{hp:eq:finite-linear}
 W(D)=W_\chi(D)\oplus W_\beta(D)\oplus W_\gamma(D)
 =W^{\leq1,\bullet}(D)[3]\oplus W^{3,\bullet}(D)[3]
                              \oplus W^{0,\bullet}(D)[1].
\end{equation}
Compact support here means compact support of the base coefficients.
Since $f$ is proper, this is equivalent to compact total-space
support for invariant forms. Haar averaging of a compact support
stays inside the inverse image of its compact projection to the
base. Thus there are support-preserving maps
\begin{equation}\label{hp:eq:IR}
 I:W(D)\longrightarrow V_H(X_D),\qquad
 R:V_H(X_D)\longrightarrow W(D),\qquad RI=1.
\end{equation}
Holomorphic translations commute with $\partial$ and
$\bar\partial$, so $I$ and $R$ commute with
\eqref{hp:eq:differential}. Write $D_W$ for the restricted
differential on $W$.

This differential is explicit. Define
\[
 C=\frac{d\Omega}{dz}\mathsf Y^{-1},\qquad
 \widetilde C=\frac{d\bar\Omega}{d\bar z}\mathsf Y^{-1}.
\]
The tilde is part of the notation: $\widetilde C$ is not the
complex conjugate of $C$. From
$dy=\mathsf Y^{-1}(\theta-\bar\theta)$ one obtains
\begin{equation}\label{hp:eq:four-blocks}
\begin{aligned}
 \partial\theta&=C\,dz\wedge\theta,&
 \bar\partial\theta&=-C\,dz\wedge\bar\theta,\\
 \partial\bar\theta&=\widetilde C\,d\bar z\wedge\theta,&
 \bar\partial\bar\theta&=-\widetilde C\,d\bar z\wedge\bar\theta.
\end{aligned}
\end{equation}
Together with differentiation of the base coefficients and the
Leibniz rule, these four blocks determine the operators on all
wedge monomials. Apply $\partial$ as an internal differential only
in $W_\chi$, with the same truncation as before.

The second block is the cotangent extension of the varying family.
Indeed, the relative holomorphic frame is the class of $\theta$
modulo base one-forms. In the flat Betti frame $dx,dy$, its
derivative is $\Omega'\,dy$. Projection modulo the relative
$(1,0)$ subspace gives the Kodaira--Spencer map
\begin{equation}\label{hp:eq:KS}
 \operatorname{KS}([\theta])=-C[\bar\theta]\otimes dz.
\end{equation}
In particular the derivatives of $\mu$ and $\sigma$ in
\eqref{hp:eq:native-periods} are retained. The finite complex
contains more than the fiberwise cohomology dimensions.

\subsection{Fiber integration and the finite comparison}
\label{hp:subsec:comparison}

The finite model inherits both brackets by actual oriented
integration along the fibers:
\begin{equation}\label{hp:eq:finite-brackets}
\begin{aligned}
 B_{W,\chi}(a,a')&=\int_D f_*
                   \bigl(Ia\wedge\partial Ia'\bigr),\\
 B_{W,\beta\gamma}(b,c)&=\int_D f_*
                   \bigl(Ib\wedge Ic\bigr).
\end{aligned}
\end{equation}
The complex orientation fixes $f_*$. Haar measure was normalized to
construct a projection; it is not substituted for the trace in
these formulas. In the local coframe, the relative trace is
determined by the actual number
\[
 \int_{F_z}\theta_1\wedge\theta_2
                    \wedge\bar\theta_1\wedge\bar\theta_2,
 \qquad F_z=f^{-1}(z),
\]
with the same orientation and wedge signs throughout.

The coherent coefficients can be identified without discarding
the differential. Define independently
\[
 E_0=\mathrm Rf_*\mathcal O_X,\qquad
 E_1=\mathrm Rf_*\Omega^1_{X/U},\qquad
 \omega_{X/U}=\Omega^2_{X/U}.
\]
Invariant relative Dolbeault forms represent these derived direct
images: this follows from the relative Fourier contraction proved
below, or from its vertical part alone. Their base Dolbeault
operators are induced by \eqref{hp:eq:four-blocks}. The absolute
cotangent complex has the additional block \eqref{hp:eq:KS}.
Wedge product and the same oriented integral give relative Serre
duality in the form
\begin{equation}\label{hp:eq:Serre}
 \mathrm Rf_*\omega_{X/U}[2]\xrightarrow{\ \cong\ }E_0^\vee.
\end{equation}
The shift $[2]$ is the complex dimension of the fiber. Explicitly,
the class of $\bar\theta_I$ pairs with
$\theta_1\wedge\theta_2\wedge\bar\theta_J$ when $I,J$ are
complementary subsets of $\{1,2\}$. The resulting matrix is
nondegenerate, with entries given by the displayed oriented trace
and its wedge signs. This proves the local identification in
\eqref{hp:eq:Serre}; holomorphic pullback and integration make it
compatible with changes of chart. Likewise
\[
 E_1\otimes E_1\longrightarrow
 \mathrm Rf_*\omega_{X/U}\longrightarrow\mathcal O_U[-2]
\]
is wedge followed by trace. On $E_1[1]$ it is the degree-zero
graded symmetric pairing used for the reduced chiral boson.
These are the coefficient pairings underlying the two blocks
in \eqref{hp:eq:finite-brackets}.

Put $\mathcal O_W(D)=(\mathsf S(W(D)),D_W)$, with those brackets.
The pushforward observable algebra is, by definition,
$(f_*\mathcal O_H)(D)=\mathcal O_H(f^{-1}D)$.

\begin{theorem}[Smooth-base Hamiltonian comparison]
\label{hp:thm:comparison}
For the selected smooth proper torus family and the resolved
observable model \eqref{hp:eq:observables}, the inclusion of invariant
total forms induces a strict unital $P_0$ factorization
quasi-isomorphism
\begin{equation}\label{hp:eq:main}
 \mathsf S(I):\mathcal O_W\longrightarrow f_*\mathcal O_H
\end{equation}
on $U$. This holds separately in either fixed symmetric-tensor
realization described above. The map $\mathsf S(R)$ is a chain
homotopy inverse. The objects, differentials and inclusion are
intrinsic under the actual affine changes of chart; no Poisson
property of $\mathsf S(R)$ is asserted.
\end{theorem}

\begin{proof}
We first construct a continuous contraction of the linear
complexes. Choose a smooth invariant Hermitian metric in the
vertical tangent directions and an invariant complex-linear
horizontal splitting. Such choices exist globally: average local
choices over the compact translation tori, then combine them with
a partition of unity on the base. They are auxiliary to the
contraction, not to $W,D_W,I$ or $R$.

In a marked chart, expand coefficients in the characters
\[
 e_k(x,y)=\exp\bigl(2\pi i(r^t x+s^t y)\bigr),
 \qquad k=(r,s)\in\mathbb Z^2\oplus\mathbb Z^2.
\]
Let $d_0$ be vertical Dolbeault differentiation. On the $k$th
character it is exterior multiplication by
\begin{equation}\label{hp:eq:symbol}
 q_k=2\pi i(r^t\Omega-s^t)\mathsf Y^{-1}\bar\theta.
\end{equation}
For $k\ne0$, this covector is nonzero. Indeed,
$r^t\Omega=s^t$ with $r,s$ real implies
$r^t\operatorname{Im}\Omega=0$, hence $r=s=0$.
Contraction by the Hermitian dual of $q_k$, divided by
$\lVert q_k\rVert^2$, gives an operator $h_0$ of degree $-1$.
Set $h_0=0$ on the zero character. With $\Pi=IR$, exterior
multiplication and contraction give
\begin{equation}\label{hp:eq:vertical-contraction}
 d_0h_0+h_0d_0=1-\Pi,\qquad
 h_0^2=0,\qquad h_0I=0,\qquad Rh_0=0.
\end{equation}

These Fourier formulas define operators on smooth compactly
supported forms. On any compact base set $K$, the real-linear
symbol map in \eqref{hp:eq:symbol} is uniformly invertible, so
\[
 \lVert q_k(z)\rVert\geq c_K|k|\qquad(z\in K, k\ne0)
\]
for some $c_K>0$. The coefficients of $h_0$ and all their base
derivatives have at most polynomial growth in $k$. Smooth Fourier
coefficients decay faster than any power, so the transformed
series and all its derivatives converge. These estimates also
give continuity on each fixed compact-support Fr\'echet space.
The construction does not enlarge the projection of the support
to the base. Properness then gives compact total support and
continuity on its inductive limit. In particular there is no
small-divisor problem: the symbol is elliptic on every nonzero
real lattice covector.

Write $D_H=d_0+\delta$. The remaining Dolbeault terms raise base
form degree. In the $\chi$ summand the additional truncated
$\partial$ raises total holomorphic degree. Give a form weight
equal to its base form degree plus its total holomorphic degree
in $V_\chi$. The latter degree is at most one and base form
degree is at most two. Thus $\delta h_0$ raises a bounded weight,
and
\[
 (\delta h_0)^4=0.
\]
For $V_\beta,V_\gamma$, the base-degree filtration alone suffices
and gives a shorter bound. Define the finite operator
\begin{equation}\label{hp:eq:homotopy}
 H=h_0(1+\delta h_0)^{-1}
   =h_0\sum_{j=0}^{3}(-\delta h_0)^j.
\end{equation}
The perturbation preserves each Fourier character. Therefore
$h_0\delta I=0$ and $R\delta h_0=0$, so the inclusion and
projection require no correction. Expanding the finite sum,
using \eqref{hp:eq:vertical-contraction} and
$(d_0+\delta)^2=0$, cancels consecutive powers of $\delta h_0$
and gives
\begin{equation}\label{hp:eq:contraction}
 D_HH+HD_H=1-IR,
 \qquad RI=1,
 \qquad D_W=RD_HI.
\end{equation}
Every term of $H$ is continuous and preserves base support.
The construction using the global invariant choices agrees with
these local Fourier formulas, so \eqref{hp:eq:contraction} is a
global identity. It is also compatible with extension by zero.

We next pass to all polynomial degrees. On the $n$th tensor power,
with the usual graded tensor signs, put
\[
 K_n=\sum_{j=1}^{n}\Pi^{\otimes(j-1)}\otimes H
                         \otimes1^{\otimes(n-j)}.
\]
The tensor differential satisfies
\[
 D K_n+K_nD
 =\sum_{j=1}^{n}\Pi^{\otimes(j-1)}\otimes(1-\Pi)
                         \otimes1^{\otimes(n-j)}
 =1-\Pi^{\otimes n}.
\]
Average $K_n$ over the finite symmetric group, with its Koszul
action, to descend the identity to the symmetric power. The
operators extend continuously to each completed projective
tensor power as well. Taking the same direct sum over $n$ gives
the asserted chain homotopy inverse; in degree zero the maps
are the identity on $\mathbb C$ and the homotopy is zero.

By \eqref{hp:eq:finite-brackets}, $I$ preserves both brackets
exactly. Hence its symmetric extension is strict $P_0$ and unital.
All the maps commute with extension by zero on the base, and
inverse images of disjoint opens are disjoint. They therefore
commute with the factorization operations in every arity.
Objectwise quasi-isomorphism transfers the factorization descent
of the stated resolved model to $\mathcal O_W$. Finally, an affine
holomorphic change of torus chart preserves invariant total
forms, their exterior operators and oriented integration. It
therefore preserves the objects and the inclusion in
\eqref{hp:eq:main}. The explicit changes of chart are described
in Section~\ref{hp:subsec:gluing}.
\end{proof}

The failure of a Poisson averaging map persists for a smooth
torus. Opposite nonzero Fourier characters can have a nonzero
integrated $\beta\gamma$ pairing, although both have zero Haar
average. Thus the proof uses the bracket-preserving inclusion
and the chain contraction for different purposes. It also keeps
the kernel resolution fixed: relative harmonic representatives
do not identify the kernel of every possible reduced
presymplectic operator with the resolution in
\eqref{hp:eq:kernel-recipe}.

\subsection{The constant-family calibration}
\label{hp:subsec:constant}

For a product $X_D=F\times D$, the period matrix is constant and
all four differential blocks in \eqref{hp:eq:four-blocks} vanish. This
specialization makes a distinction between fields and observables
visible without any variation of Hodge structure.
Write $H^{p,q}(F)=H^q(F,\Omega_F^p)$. Translation-invariant
Dolbeault forms represent these spaces. The $\chi$ generators
split as
\begin{equation}\label{hp:eq:constant-split}
\begin{aligned}
 W_\chi(D)&=Z(D)\oplus C_\chi(D),\\
 Z(D)&=\bigl(\Omega^{\mathrm{dR}}_{D,c}
                         \otimes H^{0,\bullet}(F)\bigr)[3],\\
 C_\chi(D)&=\bigl(\Omega^{0,\bullet}_{D,c}
                         \otimes H^{1,\bullet}(F)\bigr)[3].
\end{aligned}
\end{equation}
In the last line $H^{1,q}(F)$ retains its holomorphic degree one
in the total $\chi$ degree. The differential on $Z$ is the base
de Rham differential, and the differential on $C_\chi$ is
$\bar\partial_D$.

\begin{proposition}[The central polynomial factor]
\label{hp:prop:central}
For the constant family, $Z(D)$ has zero bracket with all linear
generators. Consequently the resolved observable algebra has a
central factor $\mathsf S(Z(D))$:
\begin{equation}\label{hp:eq:central-factor}
 \mathcal O_W(D)=\mathsf S(Z(D))\otimes
                 \mathsf S\bigl(C_\chi(D)\oplus
                                W_\beta(D)\oplus W_\gamma(D)\bigr),
\end{equation}
with the matching finite tensor convention. If $D$ is a disk,
$\mathsf S(Z(D))$ is not quasi-isomorphic to the unit algebra.
\end{proposition}

\begin{proof}
A $Z$ form has vertical holomorphic degree zero. On a constant
torus, $\partial$ of an invariant form differentiates only its
base coefficients. Pairing a $Z$ form with another invariant
$\chi$ form cannot produce vertical holomorphic degree two;
when both terms contain the base holomorphic form, their wedge
vanishes instead. Thus $B_{W,\chi}(Z,W_\chi)=0$.
The other mixed brackets vanish by
\eqref{hp:eq:brackets}. The symmetric algebra of the direct sum
then gives \eqref{hp:eq:central-factor}.

Compactly supported de Rham cohomology of a disk is $\mathbb C$
in degree two and zero otherwise. Hence
\begin{equation}\label{hp:eq:central-cohomology}
 H^{q-1}(Z(D))=H^{0,q}(F),\qquad q=0,1,2.
\end{equation}
The polynomial differential preserves polynomial degree, so this
nonzero linear cohomology cannot disappear in the symmetric
algebra. This proves the final assertion.
\end{proof}

The remaining two factors have the familiar chiral and
$\beta\gamma$ descriptions. The pairing on $H^{1,\bullet}(F)[1]$
is wedge and Serre trace, and the $\beta\gamma$ coefficients are
$H^\bullet(F,\mathcal O_F)$ with their graded duals. Their
dimensions, with degrees retained, are
\[
 \bigl(h^0(\mathcal O_F),h^1(\mathcal O_F),h^2(\mathcal O_F)\bigr)
 =(1,2,1),\qquad
 \bigl(h^0(\Omega_F^1),h^1(\Omega_F^1),h^2(\Omega_F^1)\bigr)
 =(2,4,2).
\]
The central factor in \eqref{hp:eq:central-factor} is an additional
part of the actual resolved observable pushforward.

By comparison, the product compactification of the field complex
contains the zero-presymplectic summand
\begin{equation}\label{hp:eq:zero-field}
 \mathcal E_{\mathrm{rad}}
 =\bigl(\Omega_D^{\mathrm{dR}}\otimes
                    H^{0,\bullet}(F)\bigr)[2]
\end{equation}
in the presymplectic field decomposition of
\cite[\S7.1.2]{SaberiWilliamsConstraints}. If Hamiltonian
observables are constructed directly from this summand, its
kernel is all of $\mathcal E_{\mathrm{rad}}$. Use the identity
as the kernel resolution map. The linear resolved observable
complex is then a shifted density dual of the cone of the
identity, hence contractible. Its symmetric algebra has
cohomology $\mathbb C$, with the scalar unit preserved. This is
also the order-zero radical case of the Hamiltonian construction.

Thus pushing forward the specified resolved observable model and
applying Hamiltonian observables to the reduced zero-pairing
field summand give different answers to this calibration:
$\mathsf S(Z)$ and the unit, respectively. The augmentation
removing $\mathsf S(Z)$ is a unital $P_0$ map, but
\eqref{hp:eq:central-cohomology} shows that it is not a
quasi-isomorphism. Identifying the two procedures would require
an additional comparison of kernel resolutions or an explicitly
specified quotient. A field-level equivalence alone supplies
neither.

\subsection{A nonzero Kodaira--Spencer coefficient at a regular value}
\label{hp:subsec:variation}

The actual elliptic equation gives a direct example in which the
mixed block of \eqref{hp:eq:four-blocks} is indispensable. At
$t=2$, equation~\eqref{hol:eq:surface} is
\[
 y^2=x^3-24x+32,
\]
with ordered real roots
\[
 e_1=4,\qquad e_2=-2+2\sqrt3,\qquad e_3=-2-2\sqrt3.
\]
Define the elliptic parameter and complete integral by
\begin{equation}\label{hp:eq:elliptic-parameter}
 m=\frac{e_2-e_3}{e_1-e_3},\qquad
 K(m)=\int_0^{\pi/2}\frac{d\phi}{\sqrt{1-m\sin^2\phi}}.
\end{equation}
Here $m$ is the square of the modulus. In the positive rectangular
marking the actual periods of $dx/(2y)$ are
\begin{equation}\label{hp:eq:elliptic-periods}
 A=\frac{2K(m)}{\sqrt{e_1-e_3}},\qquad
 B=\frac{2iK(1-m)}{\sqrt{e_1-e_3}},\qquad
 \tau=\frac BA=i\frac{K(1-m)}{K(m)}.
\end{equation}
These are the rectangular period integrals
\cite[\S23.6(iv)]{DLMFElliptic}; the chosen cycles continue
holomorphically near $t=2$.

For the cubic $x^3+A_0(t)x+B_0(t)$, put
$A_0=-3t^3(t-1)$ and $B_0=2t^4(t-1)^2$.
Since $A_0'(2)=-60$ and $B_0'(2)=128$, implicit differentiation
gives
\[
 e_i'(2)=-\frac{A_0'(2)e_i+B_0'(2)}{3e_i^2+A_0(2)}.
\]
Thus
\begin{equation}\label{hp:eq:root-jets}
\begin{gathered}
 e_1'=\frac{14}{3},\qquad
 e_2'=\frac{-7+8\sqrt3}{3},\qquad
 e_3'=\frac{-7-8\sqrt3}{3},\\
 m_*:=m(2)=\sqrt3-1,\qquad
 m_*'=\frac{2\sqrt3-3}{6}>0.
\end{gathered}
\end{equation}
The derivative of $K$ and Legendre's relation, expressed in the
parameter $m$, are
\[
 \frac{dK}{dm}=\frac{E(m)}{2m(1-m)}-\frac{K(m)}{2m},
 \qquad
 E(m)K(1-m)+E(1-m)K(m)-K(m)K(1-m)=\frac\pi2,
\]
where $E(m)=\int_0^{\pi/2}\sqrt{1-m\sin^2\phi}\,d\phi$
\cite[\S\S19.4,19.7]{DLMFElliptic}. Differentiating
\eqref{hp:eq:elliptic-periods} therefore yields
\begin{equation}\label{hp:eq:tau-derivative}
 a:=\tau'(2)
 =-\frac{i\pi m_*'}{4m_*(1-m_*)K(m_*)^2}\ne0.
\end{equation}
The factor $4$ here includes the conversion from modulus to
parameter.

Normalize $\mu$ by the $A$ period of $dx/(2y)$ and the actual
section $P=(0,\sqrt2\,t^2(t-1))$. Retain $\sigma$, the logarithm
of the original quotient-lift multiplier. Set $b=\mu'(2)$ and $c=\sigma'(2)$;
no nonvanishing assumption on $b$ or $c$ is needed. Then
\[
 \Omega'(2)=\begin{pmatrix}a&-6b\\b&c\end{pmatrix},
 \qquad \mathsf Y_* =\Omega(2)-\bar\Omega(2).
\]
Suppose only the finite-model term
$\bar\partial\theta=-C\,dt\wedge\bar\theta$ is discarded,
with the inclusion, the source differential, all periods and
every other block fixed. On $\theta_1$, the coefficient of
$dt\wedge\bar\theta$ in $D_HI-I\widetilde D_W$ is then
\begin{equation}\label{hp:eq:KS-residual}
 -[a,-6b]\mathsf Y_*^{-1}.
\end{equation}
For $\mathbf e_1=(1,0)^t$, its value on
$\mathsf Y_*\mathbf e_1/a$ is exactly $-1$. In particular the
unchanged inclusion is no longer a chain map. Multiplying
$\theta_1$ by a smooth base bump equal to one near $2$ places
this example in the actual compactly supported generator
complex. The bump-derivative terms agree on the two sides and
cancel. The failure therefore concerns the compact observable
comparison itself, not merely a choice of period frame or a
flatness equation for another connection.

\subsection{Affine overlaps, root coefficients and cusp transport}
\label{hp:subsec:gluing}

The finite model uses absolute forms, so an affine change of
origin contributes derivatives in base directions. Choose a
holomorphic lift $b(z)$ of a translation section, and write uniquely
$b=r+\Omega s$ with $r,s$ real smooth vectors. Translation
$T_b:w\mapsto w+b$ acts by
\begin{equation}\label{hp:eq:affine-pullback}
 T_b^*\theta=\theta+db-d\Omega\,s,
 \qquad
 T_b^*\bar\theta=\bar\theta+d\bar b-d\bar\Omega\,s.
\end{equation}
Changing the lift to $b+n+\Omega m$, with integral constant
$n,m$, changes $s$ to $s+m$ and leaves these formulas unchanged.
Their inverses are the pullbacks for $T_{-b}$. Integral changes
of marking act by their actual holomorphic pullbacks on the
same total-form complex. Such pullbacks preserve $W$, commute
with $I,R,\partial,\bar\partial$, and preserve the integration
pairings. Their compositions are the compositions of the
original maps, so no separate choice of homotopy enters the
overlap identities.

For the finite collars these maps can be written using the
specified integral periods. Set
\[
 (d_0,d_1)=(3,4),\qquad v_0=e_4,\qquad v_1=-f_4,
 \qquad h_i(s)=\frac{\log s}{2\pi i}v_i(s).
\]
If $\rho_i$ is the original affine deck map and
$\widetilde\rho_i$ is obtained by adding $v_i/d_i$, the actual
period identity gives
\[
 h_i(\zeta_i s)=A_i(s)h_i(s)+v_i(\zeta_i s)/d_i,
 \qquad
 T_{h_i}\rho_i=\widetilde\rho_iT_{h_i},
 \qquad \zeta_i=e^{2\pi i/d_i}.
\]
Both sides contain the original affine term. Formula
\eqref{hp:eq:affine-pullback} applies to these logarithmic
translations, with all derivatives retained. The conormal
comparisons add the windings $-2\delta_0$ and $-\delta_1$,
so the whole collar vectors remain
\[
 w_0=e_4-2\delta_0,\qquad w_1=-f_4-\delta_1,
\]
as in \eqref{hol:eq:boundary-vectors}. In particular the radial
part of $\log s$ is part of the overlap map.

On a punctured cusp disk the full periods are
\begin{equation}\label{hp:eq:cusp-periods}
 \Omega_\infty(q)=\frac{\log q}{2\pi i}I_2+B(q),
 \qquad
 B(q)=\frac1{2\pi i}
 \begin{pmatrix}\log c_1&\log c_3\\\log c_2&\log c_4\end{pmatrix},
\end{equation}
where all four $c_j$ are the actual holomorphic units in
\eqref{c26:eq:periods}. For the original divisor-line quotient
the multiplicative generators are
$(q c_1,p)$ and $(p^{-6},u q b)$ from
\eqref{hol:eq:cusp-generators}; the adapted marking carries
these units through its stated monomial change. In the notation
of \eqref{cmp:eq:cusp}, the overlap itself has the form
\[
 (\eta_1,\eta_2)\longmapsto
 \left(s_1s_2^2e^{2\pi i(\eta_1+2\eta_2)},
                     s_2e^{2\pi i\eta_2}\right),
\]
where $s_1,s_2$ denote the two unit coordinates of the retained
section. Its monomial matrix is unimodular, and division by
these units followed by the inverse monomial map gives its
inverse. Their logarithmic derivatives are included when it
acts on total forms.

The same marking determines integral cohomology transport.
For positive continuation around the cusp, in the ordered
period basis $(\alpha_1,\alpha_2,\beta_1,\beta_2)$, put
\begin{equation}\label{hp:eq:cusp-matrices}
 T=\begin{pmatrix}1&0&1&0\\0&1&0&1\\0&0&1&0\\0&0&0&1\end{pmatrix},
 \qquad
 L=\begin{pmatrix}1&0&0&0\\0&-1&0&0\\-1&1&1&0\\0&0&0&1\end{pmatrix}.
\end{equation}
The columns of $L$ express that basis in the common $e$ marking.
Equation~\eqref{c26:eq:nativemonodromy} gives
\begin{equation}\label{hp:eq:dual-transport}
 LTL^{-1}=T_\infty=
 \begin{pmatrix}2&1&1&0\\0&1&0&-1\\-1&-1&0&1\\0&0&0&1\end{pmatrix},
 \qquad
 T_{H^1}=T_\infty^{-T}.
\end{equation}
Indeed, if a cycle is transported by $T_\infty$, preservation
of its evaluation against a covector requires
$T_\infty^tT_{H^1}=I$. Higher cohomology uses the corresponding
exterior powers. Substituting $T_\infty$ itself for $T_{H^1}$
fails this identity: the $(1,1)$ entry of
$T_\infty^tT_\infty-I$ is $4$.
Thus recombining the invariant Dolbeault forms into de Rham
cohomology retains the full integral unipotent transport. The
finite discriminant representation cannot replace it.

Over a good-reduction root disk, Theorem~\ref{hp:thm:comparison}
and Proposition~\ref{hp:prop:finite} combine. The cyclic action
preserves invariant total forms. Averaging the chain homotopy
over this finite action gives an equivariant contraction, so
restriction to invariant linear sections remains a
quasi-isomorphism. The descended bracket is the cover bracket
divided by $3$ or $4$, respectively.
The resulting coefficients on the coarse disk are the invariant
pushforward of the smooth root coefficient sheaf. They need
not form an ordinary smooth vector bundle in the coarse
coordinate. For example, $|s|^2$ is invariant under root
rotation, but its expression $|t|^{2/m}$ for $t=s^m$, with
$m=3,4$, is not smooth at zero. The root coefficient sheaf,
the affine action and the degree-normalized trace are all part
of this extension.

\subsection{The remaining classical and quantum questions}
\label{hp:subsec:scope}

The construction proves a comparison of the specified resolved
classical observables on the smooth base, and finite descent at
the two root quotients whose total-space actions are free. At the toric cusp,
\eqref{hp:eq:cusp-periods} and \eqref{hp:eq:dual-transport}
describe the punctured comparison. Extending it across the
central fiber requires a Hamiltonian complex on the actual
singular fiber and its ambient neighborhood, a comparison with
the normalization and conductor, and compatibility of their
support and trace maps. Smooth-fiber Fourier contraction does
not construct these data.

The coefficients above are complex, perturbative and
polynomial. Compact lattice fields, Wu data and differential
cocycles require their own global construction. Quantization
must then supply the quantum operations and their anomalies;
states and a vertex algebra require additional structures
beyond a classical $P_0$ factorization algebra. An interacting
ADE theory is a further problem. The bosonic target-space
chiral differential operator gerbe of
Section~\ref{sec:chiral-operators} also concerns a different
observable construction. No identification of that gerbe, or
of its possible quantum realizations, with
\eqref{hp:eq:main} follows from the classical comparison.
